\sect{Теоретическое предложение: гибридная схема трекер плюс
RL-anticipation}{proposal}
%
Опираясь на проведённый разбор, мы можем теперь сформулировать
гибридную архитектуру, в которой классический вероятностный
трекер (HMM, HSMM или их гибрид) сохраняется в качестве ядра
системы, а RL-агент действует как тонкий \textit{опережающий
слой}, предсказывающий ожидаемый темп исполнения на горизонте
$200$--$500$\,мс. Подчеркнём ещё раз: всё, что излагается ниже,
имеет сугубо теоретический характер. Численных экспериментов
здесь нет, и любые гипотетические свойства предложенной схемы
подлежат проверке в последующих исследованиях.

\paragraph*{Мотивация постановки.} Из разбора в~\Secref{rl-music}
следует, что чистая RL-замена классического
трекера~\cite{Dorfer2018RL} имеет два связанных недостатка
в нашем целевом сценарии (виртуальный оркестр, концертные
произведения): требует тщательно размеченных данных уровня
\textit{позиция-в-каждый-момент}, которые крайне трудно
получить для произведений с длительными эпизодами свободного
rubato; и не использует имеющуюся \textit{априорную} структуру
партитуры (повторы, скачки, явные длительности нот),
естественно описываемую полу-марковскими
моделями~\cite{Cont2010,Nakamura2015}. С другой стороны,
классические трекеры~\cite{Cont2010,Raphael2010,Nakamura2015}
обладают зрелой математической базой и устойчивостью, но их
\textit{anticipation}-механизм~--- предсказание ожидаемого
момента следующего события~--- зашит в виде явных байесовских
обновлений, неадаптивных к индивидуальному стилю исполнителя.

\paragraph*{Архитектурная идея.} Мы рассматриваем композицию
двух модулей. Первый~--- классический вероятностный
\textit{трекер позиции}, реализованный через гибридную
полу-марковскую модель в духе~\cite{Cont2010} с явными
длительностями для нот и пауз и геометрическими длительностями
для атемпоральных объектов~\eqref{eq:cont-implicit-duration};
выходом трекера на каждом тике~$t$ являются нормированное
forward-распределение $\tilde\alpha_t \in \mathbb{R}^N$ (см.\
\Eqref{forward-rec}) и оценка наиболее вероятной позиции
$\hat\varphi(t) = \argmaxop_i \tilde\alpha_t(i)$. Второй~---
RL-агент \textit{предсказания темпа}, на вход которого
подаётся компактное представление текущего состояния трекера и
истории исполнения, а на выход~--- скалярный темповый
коэффициент~$a_t$, применяемый блоком воспроизведения на
ближайшем горизонте $H_{\text{rl}}$ (типично
$200$--$500$\,мс).

В этой постановке трекер и RL-агент работают параллельно с
одинаковой тактовой частотой; информационный поток~---
односторонний (от трекера к агенту), что упрощает теоретический
анализ устойчивости по сравнению с полной coupled-схемой
Конта~\cite{Cont2010}. Графически архитектура изображена на
\Figref{rl-agent}.

\par\addvspace{2pt}
{\centering
\refstepcounter{figure}\figlab{rl-agent}%
\resizebox{\columnwidth}{!}{% =====================================================================
% TikZ figure: RL anticipation agent (BLACK & WHITE)
% Compact layout: 3 boxes top, 2 boxes bottom; fits into one column.
% Used in: sections/07-rl-anticipation.tex
% =====================================================================
\begin{tikzpicture}[
    node distance = 4mm and 3mm,
    every node/.append style={font=\scriptsize},
    block/.style ={rectangle, draw=black, thick, rounded corners=2pt,
                   minimum height=9mm, minimum width=20mm,
                   align=center, font=\scriptsize},
    arrow/.style ={-Stealth, thick}]

  \node[block]                  (track)  {Базовый\\трекер};
  \node[block, right=of track]  (state)  {Состояние\\$s_t$};
  \node[block, right=of state]  (policy) {Политика\\$\pi_\theta$};

  \node[block, below=7mm of policy] (env)    {Среда\\(рендерер)};
  \node[block, left=of env]         (reward) {Награда\\$r_t$};

  \draw[arrow] (track)  -- (state);
  \draw[arrow] (state)  -- (policy);
  \draw[arrow] (policy) -- node[right, font=\tiny] {$a_t$} (env);
  \draw[arrow] (env)    -- (reward);
  \draw[arrow] (reward.west) -| (track.south);

\end{tikzpicture}
}\par
\vspace{2pt}
\footnotesize
\figurename~\thefigure. Схема предлагаемой гибридной
архитектуры. Базовый вероятностный трекер выдаёт состояние
$s_t$, RL-агент~--- темповый коэффициент $a_t \in[0{,}5;\,2{,}0]$
на горизонте $200$--$500$\,мс, среда (рендерер плюс симулятор
солиста) формирует вознаграждение $r_t$, которое замыкает
обратную связь в обучении.\par}
\addvspace{4pt}

\paragraph*{Формальная постановка MDP.} Пусть $\pi_0\in\mathbb{R}_+$~---
номинальный темп партитуры в beat/секунду. RL-агент решает MDP
$\mathcal{M}^{\text{tempo}} = (\mathcal{S}^{\text{tempo}},
\mathcal{A}^{\text{tempo}}, \mathcal{T}^{\text{tempo}}, r,\gamma)$
с тиком, совпадающим с тиком трекера. Состояние
\begin{equation}
s_t = \bigl(\,\hat\alpha_t,\;\tau_{t-K:t},\;
e_{t-K:t},\;\hat\varphi(t)/N\,\bigr),
\eqlab{rl-state-proposal}
\end{equation}
где $\hat\alpha_t\in\mathbb{R}^L$~--- сниженная по размерности
версия forward-распределения (например, через DCT-проекцию,
сохраняющую первые $L$~коэффициентов), $\tau_{t-K:t}\in
\mathbb{R}^K$~--- последние $K$ оценок темпа (нормированных к
$\pi_0$), $e_{t-K:t}\in\mathbb{R}^K$~--- последние $K$ значений
эмиссионной ошибки $-\log b_{\hat\varphi}(o_t)$,
$\hat\varphi(t)/N\in[0,1]$~--- относительное положение в
партитуре. Действие $a_t\in\mathcal{A}^{\text{tempo}} =
[a_{\min}, a_{\max}]$~--- одномерный мультипликативный темповый
коэффициент, применяемый рендерером в окне следующих
$H_{\text{rl}}$ миллисекунд. Дисконт $\gamma=0.95$ соответствует
эффективному горизонту $\sim H_{\text{rl}}$ при тике $10$\,мс.

\paragraph*{Структура функции вознаграждения.} В предлагаемой
постановке функция вознаграждения сочетает три компонента,
каждый из которых имеет ясную музыкальную интерпретацию:
рассогласование между моментом отрисовки оркестрового сэмпла и
фактическим onset солиста, потеря выравнивания базового
трекера и штраф за резкие изменения темпа. Формально
\begin{equation}
\begin{aligned}
r_t = \,&-\,\bigl|\,t^{\text{render}}_t - t^{\text{perf}}_t\,\bigr|
       \,-\, \lambda\,\mathcal{L}_{\text{align}}(t) \\
       &-\, \mu\,(a_t - a_{t-1})^2,
\end{aligned}
\eqlab{rl-reward-proposal}
\end{equation}
с гиперпараметрами $\lambda,\mu>0$. Первая компонента
непосредственно отражает воспринимаемое рассинхронирование;
вторая~--- унаследованное качество выравнивания базового
трекера; третья~--- гладкость управления, акустически
обеспечивающая отсутствие «дёрганья» оркестра. Такая структура
функции вознаграждения сходна с предложенной в~\cite{Dorfer2018RL}
по формальной структуре, но отличается тем, что не сводится к
геометрической дистанции в партитуре~---
вознаграждение здесь определено непосредственно через
\textit{временное} рассогласование рендерера и солиста и
\textit{темповую} гладкость.

\paragraph*{Алгоритмическая схема.} В первой стадии
обучения (behavioural cloning) учительская траектория темпа
$a^*_t$ извлекается из выровненной MIDI-разметки как сглаженная
производная времени; политика $\pi_\theta$ инициализируется
минимизацией~$\mathcal{L}_{\text{BC}}$ (см.\ \Eqref{bc-loss}).
Во второй стадии (PPO fine-tuning) политика дообучается
с использованием~$L^{\text{CLIP}}$ из~\Eqref{ppo-clip}; в
качестве оценки преимуществ применяется GAE~\Eqref{gae-def} с
$\lambda_{\text{GAE}}=0.95$ и $\gamma=0.95$. Архитектура
политики~--- многослойный перцептрон $\mathbb{R}^d \to 128 \to
64 \to 1$, где $d$~--- размерность~$s_t$ из~\Eqref{rl-state-proposal};
для учёта временного контекста дополнительно может тестироваться
GRU-вариант. Полное описание алгоритма с учётом гиперпараметров
вынесено в~\Secref{appendix-rl-detailed}.

\paragraph*{Положение относительно существующих работ.} В
рамках классификации обзора~\Secref{rl-music} предлагаемая схема
не попадает ни в одну из четырёх существующих групп. От группы
(i)~\cite{Dorfer2018RL} она отличается тем, что классический
трекер сохраняется как ядро и не заменяется RL-агентом. От
группы (ii)~\cite{Peter2023}~--- тем, что RL-агент решает задачу
\textit{предсказания темпа} (непрерывная задача регрессии в
будущее), а не выбора кандидатного onset’а в окне (дискретная
задача классификации в текущем моменте), и работает не на
символьном MIDI, а в общей постановке. От группы
(iii)~\cite{Jiang2020RLDuet,Wu2024ReaLchords}~--- тем, что
музыкальный материал не \textit{генерируется} агентом, а
заранее задан партитурой и лишь воспроизводится с
адаптивным темпом. От группы (iv)~\cite{Jaques2017SeqTutor}~---
тем, что предобученной модели здесь нет: политика инициализируется
behavioural cloning по выровненным
данным, а не дообучается из MLE-LSTM.

\paragraph*{Ожидаемые свойства.} Предлагаемая схема обладает
рядом теоретически ожидаемых свойств, проверка которых
является направлением последующих исследований. Во-первых,
сохранение классического трекера как ядра обеспечивает
\textit{вероятностную интерпретируемость} оценки позиции через
forward-распределение $\tilde\alpha_t$. Во-вторых, отдельный
лёгкий RL-слой даёт возможность \textit{адаптации} темпового
управления под индивидуального солиста через дообучение
политики на собственных записях исполнителя без
переобучения трекера. В-третьих, разделение задач (трекер
оценивает позицию, RL-агент предсказывает темп) формально
упрощает диагностику и аблации: каждый компонент можно
отключать и оценивать вклад отдельно.

\paragraph*{Известные ограничения и риски.} Предлагаемая схема
не лишена ограничений, открыто обсуждаемых ниже. Первое~---
\textit{двойная зависимость}: качество RL-предсказания темпа
зависит от качества forward-оценки трекера (через признаки
$\hat\alpha_t$ в~\Eqref{rl-state-proposal}); следовательно,
ошибки трекера могут амплифицироваться RL-агентом. Второе~---
\textit{distributional shift} при behavioural cloning: учительская
траектория темпа извлекается из доступной разметки, а в режиме
инференса политика встречает состояния, не покрытые обучающим
распределением. Стандартный приём~--- KL-контроль
\Eqref{kl-control}~\cite{Jaques2017SeqTutor,Wu2024ReaLchords}.
Третье~--- предположение об \textit{одностороннем потоке}
информации (от трекера к агенту, без обратной связи): в полной
теоретической постановке Конта~\cite{Cont2010} темпо-агент
обновляет вероятности позиции, замыкая обратную связь;
формальное теоретическое исследование устойчивости полной
двусторонней схемы выходит за рамки настоящей работы и
обозначено как направление развития.
